\subsection{实数空间到射影空间中的嵌入}

第 2 小节的命题 5 表明，若在 \(P_{n}\) （ \(n \geqslant 0\)）中给定一个射影超平面 H ，就存在 \(R^{n}\) 到该超平面的补集 \(\complement H\) 上的一个同胚。一旦选定了 H ，用命题 5 中定义的同胚把 \(R^{n}\) 与 \(\complement H\) 等同起来常常是方便的；这时称射影超平面 H 位于“无穷远”，它的点和子集也是如此。通常取 H 为方程为 \(x_{0}=0\) 的“坐标”超平面，于是点 \(z=(z_{i})\) （\(R^{n}\) 的）就与 \(P_{n}\) 中齐次坐标为 \(1, z_{1}, z_{2}, \ldots, z_{n}\) 的点等同。

一旦作了这个等同，在 \(P_{n}\) 中，任何仿射线性簇 V （p 维的，\(R^{n}\) 中的）的闭包是一个 p 维射影线性簇，不含于无穷远超平面中，且等同于由 V 生成的射影线性簇。反之，每个射影线性簇 P （p 维的，不含于无穷远超平面的）在 \(R^{n}\) 上的迹是一个 p 维仿射线性簇，它在 \(P_{n}\) 中的闭包就是 P。

在特例 \(n = 1\) ，无穷远超平面是一个点；由于 \(\mathbf{P}_1\) 是紧的，由亚历山德罗夫定理（第 I 章，§ 9，第 8 小节，定理 4），\(\mathbf{P}_1\) 同胚于空间 \(\tilde{\mathbf{R}}\) ——由局部紧空间 \(\mathbf{R}\) 添加一个点（“无穷远点”）而紧化所得到的——。由 § 2 第 4 小节的命题 4，我们还看出实射影直线 \(\mathbf{P}_1(\mathbf{R})\) 同胚于圆周 \(\mathbf{S}_1\) 以及环面 \(\mathbf{T}\)。

另一方面，若 \(n > 1\) ，\(\mathbf{P}_n(\mathbf{R})\) 不同胚于 \(\mathbf{S}_n\) ，这一点我们以后将会看到（参见习题 4）。

空间 \(\tilde{R}\) 的“无穷远点”记作 \(\infty\) ，不带正负号。重要的是不要把 \(\tilde{R}\) 与第 IV 章，§ 4 中定义的扩充实直线 \(\overline{R}\) 相混淆，后者有两个“无穷远点”；事实上，\(\tilde{R}\) 同胚于 \(\overline{R}\) 的商空间——把两点 \(+\infty\) 与 \(-\infty\) 等同起来所得到的。

